% Evidence-only paint-order diagnostic for FIG-P634-01.
% This is a byte-for-byte semantic/layout copy of the assigned source except
% that the four foreground halo nodes are fully transparent.  It reveals the
% hatch foreground that exists before the real opaque halo is painted.
\tikzset{
  slfig-FIG-P634-01/.style={font=\fontsize{9.6pt}{11.5pt}\selectfont,every path/.append style={line cap=round,line join=round}},
  sl634-old/.style={draw=SLRuleGray,densely dotted,fill=white,minimum width=13.2mm,minimum height=9.5mm,line width=.78pt,inner sep=2pt},
  sl634-done/.style={draw=SLBlue,fill=SLBlue!4,pattern=north east lines,pattern color=SLRuleGray,minimum width=13.2mm,minimum height=9.5mm,line width=.95pt,inner sep=2pt},
  sl634-halo/.style={draw=none,fill opacity=0,text opacity=0,inner xsep=1.5pt,inner ysep=1.25pt},
  sl634-now/.style={draw=SLGold,fill=SLGold!9,minimum width=13.2mm,minimum height=9.5mm,line width=1.05pt,inner sep=2pt},
  sl634-card/.style={draw=SLRuleGray,rounded corners=1.8pt,fill=white,line width=.55pt,inner xsep=5pt}
}
\pgfkeysifdefined{/tikz/alt}{}{\tikzset{alt/.code={}}}
\begin{figure}[htbp]
\centering
\begin{tikzpicture}[slfig-FIG-P634-01,>=Stealth,every node/.style={font=\fontsize{9.6pt}{11.5pt}\selectfont,align=center},
  alt={对象—关系—结论：固定顺序为 1,2,…,j−1,j,j+1,…,d；第 j 步的 x^[j] 左侧是同轮新值且右侧是上一轮旧值；只有 x^[d]=x^(t) 是轮末样本}]
  \node[font=\fontsize{10.6pt}{12.7pt}\selectfont\bfseries,text=SLBlue] at (0,2.16) {一轮系统扫描的坐标带};
  \node at (-5.2,1.48) {$1$};
  \node at (-3.7,1.48) {$2$};
  \node at (-2.2,1.48) {省略};
  \node at (-.7,1.48) {$j$ 前一位};
  \node at (.8,1.48) {$j$};
  \node at (2.3,1.48) {$j$ 后一位};
  \node at (3.8,1.48) {省略};
  \node at (5.3,1.48) {$d$};
  \draw[-{Stealth[length=1.8mm,width=1.25mm]},draw=SLGray,line width=.60pt] (-5.52,1.10)--(5.62,1.10)
    node[right=4pt,font=\fontsize{9.6pt}{11.5pt}\selectfont] {更新顺序};
  \node[sl634-done] at (-5.2,0) {};
  \node[sl634-done] at (-3.7,0) {};
  \node[sl634-done] at (-2.2,0) {};
  \node[sl634-done] at (-.7,0) {};
  \node[sl634-halo] at (-5.2,0) {坐标\\$1$};
  \node[sl634-halo] at (-3.7,0) {坐标\\$2$};
  \node[sl634-halo] at (-2.2,0) {中间\\坐标};
  \node[sl634-halo] at (-.7,0) {坐标 $j$\\前一位};
  \node[sl634-now] at (.8,0) {坐标\\$j$};
  \node[sl634-old] at (2.3,0) {坐标 $j$\\后一位};
  \node[sl634-old] at (3.8,0) {中间\\坐标};
  \node[sl634-old] at (5.3,0) {坐标\\$d$};
  \node[font=\fontsize{9.6pt}{11.5pt}\selectfont,text=SLBlue] at (-2.9,-.76) {同轮新值};
  \node[font=\fontsize{9.6pt}{11.5pt}\selectfont,text=SLGold] at (.8,-.76) {本步新值};
  \node[font=\fontsize{9.6pt}{11.5pt}\selectfont,text=SLTextGray] at (3.8,-.76) {上一轮旧值};
  \node[sl634-card,minimum width=118mm,minimum height=10.5mm] at (0,-1.52) {};
  \node[font=\fontsize{10.0pt}{12.0pt}\selectfont] at (0,-1.25) {$x^{[j]}$ 轮内状态};
  \node[font=\fontsize{9.8pt}{11.7pt}\selectfont,text=SLBlue] at (-2.65,-1.81)
    {坐标 $1$ 至 $j$\quad 同轮新值};
  \node[font=\fontsize{9.8pt}{11.7pt}\selectfont,text=SLTextGray] at (2.65,-1.81)
    {坐标 $j$ 后一位至 $d$\quad 上一轮旧值};
  \node[sl634-card,minimum width=123mm,minimum height=11mm] at (0,-2.72) {};
  \node[inner sep=2.5pt,font=\fontsize{10.0pt}{12.0pt}\selectfont] (sweepend) at (-2.65,-2.87) {$x^{[d]}$};
  \node[inner sep=2.5pt,font=\fontsize{10.0pt}{12.0pt}\selectfont] (roundstate) at (.85,-2.87) {$x^{(t)}$};
  \node[inner sep=2.5pt,font=\fontsize{9.8pt}{11.7pt}\selectfont,text=SLTextGray] (sample) at (4.25,-2.87) {轮末样本};
  \draw[<->,draw=SLGray,line width=.60pt] (sweepend.east)--(roundstate.west)
    node[midway,above=4pt,font=\fontsize{9.6pt}{11.5pt}\selectfont,text=SLTextGray] {同一状态};
  \draw[-{Stealth[length=1.5mm,width=1.05mm]},draw=SLGray,line width=.60pt] (roundstate.east)--(sample.west)
    node[midway,above=4pt,font=\fontsize{9.6pt}{11.5pt}\selectfont,text=SLTextGray] {仅此记录};
\end{tikzpicture}
\captionsetup{width=.94\linewidth}
\caption[系统扫描按固定次序立即写回并区分轮内状态与轮末样本]{系统扫描按固定次序立即写回\quad 第 $j$ 步的轮内状态 $x^{[j]}$ 左侧为同轮新值而右侧为上一轮旧值\quad 完成第 $d$ 步后 $x^{[d]}$ 与 $x^{(t)}$ 才是同一状态并记作轮末样本}
\label{fig:V5-C04-coordinate-sweep-pre-occlusion}
\end{figure}
